\documentclass[../Main.tex]{subfiles}

\begin{document}

\section{Entwicklung} \label{sec:Entwicklung}
In diesem Kapitel gehen wir durch unser Prototyp Level, das wir so designt haben, um unsere Designvorschläge aus Kapitel \ref{sec:Designziele} zu überprüfen. Da sich unsere Erkenntnisfrage in Gegnerverhalten und Gegnerplatzierung aufteilen lässt und wir deshalb auch Designvorschläge so aufgeteilt haben, behandeln wir in den nächsten 2 Unterkapiteln diese 2 Themenbereiche individuell behandelt. Dabei werden wir jedem Designvorschlag konkrete Umsetzungen zuweisen und erklären, weshalb.

\subsection{Spieler-Gegner Interaktion}
Die Designvorschläge waren ja:
\begin{description}[style=nextline, leftmargin=0pt, itemsep=0.8em]
    \item[\textbf{Designziel 1: Überlegenheit als Grundzustand}]
        Der Spieler soll sich seinen Gegnern grundsätzlich überlegen fühlen, wobei dieser Zustand gezielt durchbrochen werden kann. Der Nutzen liegt im Spielgefühl: Aggression soll sich sicher und belohnt anfühlen, während Spannung erst dort entsteht, wo die Kontrolle kurzzeitig bedroht wird. Hergeleitet ist das Ziel aus dem Kontrast der beiden Referenzspiele, der dauerhaften Überlegenheit in \game{NieR: Automata} (Abschnitt~\ref{sec:ReferenzspielanalyseNier}) gegenüber dem prekären Gleichgewicht in \game{Ninja Gaiden} (Abschnitt~\ref{sec:ReferenzspielanalyseNinja}), sowie aus dem Prinzip des emotionsgeleiteten Designs (Abschnitt~\ref{sec:Pacing}).

    \item[\textbf{Designziel 2: Begrenzung gleichzeitiger Angreifer}]
        Auch bei vielen Gegnern sollen nie alle gleichzeitig angreifen, sondern nur eine begrenzte Zahl. Der Nutzen betrifft Schwierigkeit und Spielfluss, da grosse Gruppen dadurch lesbar und fair bleiben, statt den Spieler zu überfordern. Beide Referenzspiele nutzen dieses Mittel, in \game{Ninja Gaiden} mit maximal etwa vier aktiven Gegnern, in \game{NieR: Automata} mit vermutlich höchstens zwei gleichzeitig angreifenden Nahkämpfern (Abschnitt~\ref{sec:Vergleich}), und es entspricht dem Pacing-Gedanken der Recherche (Abschnitt~\ref{sec:Pacing}).

    \item[\textbf{Designziel 3: Ein taktisches Dilemma durch zwei Gegnerrollen}]
        Der Kampfabschnitt soll Nah- und Fernkämpfer kombinieren, sodass der Spieler entscheiden muss, welche Bedrohung er zuerst ausschaltet. Der Nutzen liegt in der Schwierigkeit, die dadurch taktisch statt rein mechanisch (etwa durch mehr Lebenspunkte der Gegner) entsteht und schon kleine Gruppen anspruchsvoll macht. Grundlage ist, dass genau diese Kombination in beiden Referenzspielen den taktischen Druck trägt (Frage~D2 beider Analysen), sowie das Prinzip der Gegnerrollen im Encounter-Design (Abschnitt~\ref{sec:Encounter}).

    \item[\textbf{Designziel 4: Klar telegrafierte Angriffe}]
        Bevorstehende Gegnerangriffe sollen deutlich angekündigt werden, etwa durch sichtbare Vorzeichen. Der Nutzen betrifft Schwierigkeit und Spielgefühl, weil der Spieler auch im Gruppenkampf rechtzeitig reagieren kann und Treffer als eigene Fehler statt als Willkür erlebt. Vorbild ist das deutliche Telegrafieren in \game{NieR: Automata} (Frage~M2), theoretisch gestützt durch das Prinzip der Lesbarkeit (Abschnitt~\ref{sec:Encounter}).
\end{description}

\subsubsection{Gestaltung des Kampfabschnitts}

\begin{itemize}[leftmargin=1.5cm]
    \item \textcolor{red}{[Platzhalter] Konkrete Gestaltung des Kampfabschnitts aus den Designzielen: Raumlayout, Gegnerpositionen und Spawnpunkte, Wellenabfolge sowie Höhen- und Raumvariation.}
    \item \textcolor{red}{Direkteste Grundlagen laut Wissenslandkarte: Kapitel 1 (Encounter-Design), Kapitel 3 (Raumtypen) und Kapitel 4 (Pacing).}
    \item \textcolor{red}{Planung vor der Greybox gegebenenfalls mit einem Molekül- oder Bubble-Diagramm.}
\end{itemize}

\subsubsection{Umsetzung im Prototyp}

\begin{itemize}[leftmargin=1.5cm]
    \item \textcolor{red}{[Platzhalter] Umsetzung der Gestaltung in Godot: Greybox, Szenenaufbau und Gegnerplatzierung im Level.}
    \item \textcolor{red}{Architekturbezug laut Wissenslandkarte: Kapitel 2 (Kernmechanik) und Kapitel 10 (Modularität) begründen die Godot-Entscheidungen.}
    \item \textcolor{red}{Technische Details siehe Kapitel 5.}

\end{itemize}
\end{document}
